\appendix
\section{Local stability in edit distance}\label{app:local-stability}

We give a self-contained local argument. It permits fractional edge
values and triangles, which may arise during the modifications in the
proof. Accordingly, throughout this appendix put
\[
 P(W)=\frac1{10}\int_{[0,1]^5}
       \prod_{ij\in E(C_5)}W(x_i,x_j)
       \prod_{\substack{i<j\\ij\notin E(C_5)}}(1-W(x_i,x_j))
       \,dx_1\cdots dx_5,
\]
\[
 F(W)=\frac{P(W)}{E(W)^{5/2}},\qquad
 R_{\mathrm{ind}}(W)=5^{5/2}F(W).
\]
These definitions agree with $P(W)=C(W)$ and $R_{\mathrm{ind}}(W)=R(W)$
for every graphon in the main text. Let $W_0$ be the Clebsch graphon on
a fixed partition $A_1,\ldots,A_{16}$ into sets of measure $a=1/16$.
As before, $\mathcal E_{\mathrm{Cl}}$ includes balanced Clebsch graphons
on any positive active mass, with isolates on the complement. Set
\[
 d_1(W,\mathcal E_{\mathrm{Cl}})
   =\inf_{U\in\mathcal E_{\mathrm{Cl}}}\lVert W-U\rVert_1.
\]
All measurable partitions are allowed in the infimum.

\begin{theorem}[Local quadratic stability]\label{thm:local-stability}
There exist $\delta_0>0$ and $c_{\mathrm{loc}}>0$ such that every graphon
$W:[0,1]^2\to[0,1]$ with $\lVert W-W_0\rVert_1<\delta_0$ satisfies
\begin{equation}\tag{76}\label{eq:local-stability}
 R_\ast-R_{\mathrm{ind}}(W)
   \geq c_{\mathrm{loc}}d_1(W,\mathcal E_{\mathrm{Cl}})^2.
\end{equation}
Equality in $R_{\mathrm{ind}}(W)\leq R_\ast$ holds precisely when
$W\in\mathcal E_{\mathrm{Cl}}$ up to null sets.
\end{theorem}

\subsection{Three finite properties of the Clebsch graph}

Write the five generators of $\mathrm{Cl}$ as $s_1,\ldots,s_5$, with
$\sum_i s_i=0$. Their only nonempty dependence is the full five-element
set. Relative to a vertex $u$, its neighbors are $u+s_i$ and its ten
nonneighbors, other than itself, are $u+s_i+s_j$, indexed by
$\{i,j\}\in\binom{[5]}{2}$. A pair vertex is adjacent to the two
neighbors indexed by its elements; two pair vertices are adjacent
exactly when their pairs are disjoint. In particular, the ten pair
vertices induce the Petersen graph $\mathrm{KG}(5,2)$.

We will also use the following counts, including their normalizations.
Every closed walk of length five uses each generator once, and is a
simple pentagon. Hence $\mathrm{Cl}$ has $16\cdot5!/10=192$ pentagons
and each vertex belongs to $5!/2=60$ of them. A fixed edge belongs to
$4!=24$ pentagons: orient it from its specified first endpoint and
permute the remaining four generators. For a fixed nonedge from $0$ to
$s_i+s_j$, a rooted closed walk reaches the second endpoint at step two
or at step three. There are $2!3!$ orders in each case. Dividing the
total by the two orientations gives $12$ pentagons through that pair.

\begin{lemma}[The neighborhood inequality]\label{lem:neighborhood}
Add a vertex to $\mathrm{Cl}$ with neighborhood $N\subseteq V(\mathrm{Cl})$,
and let $c(N)$ count the induced pentagons containing it. Then
\[
 c(N)\leq12|N|.
\]
Equality holds exactly when $N=\varnothing$ or $N=N_{\mathrm{Cl}}(v)$
for a vertex $v$.
\end{lemma}
\begin{proof}
Deleting the added vertex from a pentagon leaves an induced $P_4$ with
its endpoints in $N$ and internal vertices outside $N$. For $u\in N$,
let $c_u(N)$ count these paths with endpoint $u$, so that
$2c(N)=\sum_{u\in N}c_u(N)$.

Use the pair labeling relative to $u$. Let $U$ be its ten pair vertices,
$T=N\cap U$, and $D=\{i:u+s_i\in N\}$. A path counted by $c_u(N)$
has the form $u,u+s_i,b,v$, where $b\in U\setminus T$, $v\in T$, the
pairs $b,v$ are disjoint, and $i\in b\setminus D$. These conditions
are also sufficient. Thus, writing $P=\mathrm{KG}(5,2)$,
\begin{equation}\tag{77}\label{eq:neighborhood-path-count}
 c_u(N)=\sum_{b\in U\setminus T}(2-|b\cap D|)d_T(b)
   \leq2e_P(T,U\setminus T).
\end{equation}

Here the maximum cut of $P$ has size twelve. To verify this directly,
$P$ has parameters $(10,3,0,1)$: disjoint pairs have no common
neighbor, while two intersecting pairs have the unique complementary
pair as common neighbor. It has no four-cycle. For any fixed edge,
choose one of the two other neighbors at each endpoint. The two choices
are nonadjacent and have a unique common neighbor; they therefore
complete the edge to a pentagon. This gives four pentagons through
every edge, and $15\cdot4/5=12$ pentagons in total. Every cut leaves
at least one edge of each odd cycle uncut. Since an edge lies in four
pentagons, at least three of the fifteen edges are uncut. A star
\[
 F_i=\{\{i,j\}:j\ne i\}
\]
is an independent four-set with twelve edges leaving it, so the bound
is attained.

Moreover every twelve-edge cut has a side of size four, and that side
is independent. Indeed, a side of size at most three has at most nine
outgoing edges, and a side of size five has an odd number of outgoing
edges. Four pairwise intersecting two-subsets form a star: if
$\{1,2\},\{1,3\}$ are present and some pair omits $1$, that pair must
be $\{2,3\}$ and there can be no fourth distinct pair. Thus the
maximum cut sides are precisely $F_i$ and $U\setminus F_i$.

Equation (77) proves $c_u(N)\leq24$. If equality holds and $T=F_i$,
every pair outside this star has two neighbors in it. The loss term in
(77) vanishes only if $D\subseteq\{i\}$. Consequently
$N=N_{\mathrm{Cl}}(u+s_i)$ or $N=N_{\mathrm{Cl}}[u+s_i]$, of size five
or six. If $T=U\setminus F_i$, every pair in $F_i$ has three neighbors
in $T$, and the pairs in $F_i$ use all five indices. Hence
$D=\varnothing$ and $N=\{u\}\cup(U\setminus F_i)$ has size seven;
in this case $u$ has no neighbor in $N$.

Suppose $N\ne\varnothing$ and $c(N)=12|N|$. Then every $u\in N$ has
$c_u(N)=24$. The size-six case is impossible: at the center of a closed
neighborhood, the first internal vertex of a counted path cannot lie
outside $N$, so its count is zero. The size-seven case is also
impossible. The six pair vertices in that case contain an edge, since
an independent set in $\mathrm{KG}(5,2)$ has at most four vertices by
the preceding intersecting-pair argument. An endpoint of this edge has
a neighbor in $N$, whereas the equality classification at that endpoint
permits a neighbor in $N$ only in the size-six case. Therefore $N$ is
a vertex neighborhood. Conversely, adding a clone of a vertex creates
exactly its sixty pentagons. No such pentagon can contain both clones,
since $C_5$ has no pair of twins. Thus $c(N)=60=12|N|$ for a vertex
neighborhood, and the empty case is immediate.
\end{proof}

\begin{lemma}[Two-vertex coefficients]\label{lem:two-vertex}
Add two vertices with old neighborhoods $N_{\mathrm{Cl}}(i)$,
$N_{\mathrm{Cl}}(j)$, and prescribe whether they are adjacent. The
numbers of induced pentagons containing both added vertices and three
old vertices are
\begin{center}
\begin{tabular}{lrr}
\toprule
&added pair absent&added pair present\\
\midrule
$i=j$&$0$&$0$\\
$ij\in E(\mathrm{Cl})$&$0$&$24$\\
$i\ne j,\ ij\notin E(\mathrm{Cl})$&$12$&$6$\\
\bottomrule
\end{tabular}
\end{center}
\end{lemma}
\begin{proof}
For equal types the added vertices are twins, whether adjacent or not,
so cannot both occur in an induced pentagon. For distinct types, when
their prescribed adjacency agrees with that of $i,j$, the counts are
the ordinary pair counts $24$ and $12$ above. A pentagon cannot use
two vertices from the same clone class, so the old vertices of types
$i,j$ introduce no extra completions. When $ij$ is an edge but the
added pair is absent, the two added vertices have no common old
neighbor. This rules out a pentagon, in which every nonadjacent pair
has a common neighbor.

For the last entry translate $i$ to $0$ and write $j=s_1+s_2$. Their
exclusive old neighbors are, respectively,
\[
 \{s_3,s_4,s_5\},\qquad
 \{s_3+s_4,s_3+s_5,s_4+s_5\}.
\]
A pentagon containing the added edge must join one exclusive neighbor
on either side through a common old vertex outside
$N_{\mathrm{Cl}}(i)\cup N_{\mathrm{Cl}}(j)$. The two exclusive
neighbors must be nonadjacent. Thus a pair from
$\binom{\{3,4,5\}}{2}$ is matched with its unique missing generator,
giving three choices. For example $s_3+s_4$ and $s_5$ have common
neighbors $s_1+s_5,s_2+s_5$, both outside the forbidden union. There
are exactly two completions for each choice, giving six in total.
\end{proof}

\begin{lemma}[The template Hessian]\label{lem:template-hessian}
For the Clebsch vertex-weight polynomials $P(x)$ and $E(x)$,
\[
 \nabla^2\log(P/E^{5/2})(\mathbf1)=\frac{J-16I}{256}.
\]
Consequently there is $c_H>0$ such that, for class masses $s_i$
sufficiently close to $a$, with $q=\sum_i s_i$ and $T_s$ the
corresponding Clebsch step graphon,
\begin{equation}\tag{78}\label{eq:hessian-defect}
 F(W_0)-F(T_s)\geq c_H\sum_i(s_i-q/16)^2.
\end{equation}
\end{lemma}
\begin{proof}
Let $A$ be the Clebsch adjacency matrix. The counts above give
\[
 P(\mathbf1)=192,\qquad\nabla P(\mathbf1)=60\mathbf1,\qquad
 \nabla^2P(\mathbf1)=12(J-I+A),
\]
whereas $E(\mathbf1)=40$, $\nabla E(\mathbf1)=5\mathbf1$ and
$\nabla^2E=A$. Hence the logarithmic Hessian is
\[
 \frac{12(J-I+A)}{192}-\frac{3600J}{192^2}
 -\frac52\left(\frac A{40}-\frac{25J}{40^2}\right)
 =\frac{J-16I}{256}.
\]
The gradient vanishes at $\mathbf1$. On the hyperplane $\sum_i h_i=0$,
Taylor's formula gives
\[
 \log F(\mathbf1+h)
 =\log F(\mathbf1)-\frac{\lVert h\rVert_2^2}{32}
     +O(\lVert h\rVert_2^3).
\]
Apply this with $h_i=16s_i/q-1$ and use scale invariance. Since $q$
stays close to one, this gives (78) with a uniform positive constant.
\end{proof}

\subsection{Fractional row profiles}

Let $\mathcal M\subseteq\{0,1\}^{16}$ consist of the zero vector and
the sixteen Clebsch neighborhood vectors. Extend $c(N)$ multilinearly
to $p\in[0,1]^{16}$:
\[
 c(p)=\sum_{\substack{U\subseteq V(\mathrm{Cl})\\
                      \mathrm{Cl}[U]\cong P_4}}
       \prod_{\substack{i\in U\\i\text{ an endpoint}}}p_i
       \prod_{\substack{j\in U\\j\text{ internal}}}(1-p_j).
\]
This is the expectation of $c(N)$ for an independent Bernoulli subset
with inclusion probabilities $p_i$. The neighborhood lemma implies
the quantitative bound
\begin{equation}\tag{79}\label{eq:fractional-profile-distance}
 12\sum_i p_i-c(p)\geq\frac{\operatorname{dist}_1(p,\mathcal M)}{272}.
\end{equation}
For completeness, let $\epsilon$ be the probability that the Bernoulli
subset lies outside $\mathcal M$. The Boolean deficit is an integer
and at least one there, so the left side is at least $\epsilon$. No
two members of $\mathcal M$ differ in a single coordinate: the
distances from zero to a neighborhood are five, and between distinct
neighborhoods are six or ten. Conditioning on all other coordinates
therefore gives $\min(p_i,1-p_i)\leq\epsilon$ for each $i$. Round
every coordinate to its nearest Boolean value, obtaining $v$. If
$\epsilon<1/17$ and $v\notin\mathcal M$, then
\[
 \epsilon\geq\Pr(Y=v)
 \geq1-\sum_i\min(p_i,1-p_i)\geq1-16\epsilon>\epsilon,
\]
a contradiction. In this case
$\operatorname{dist}_1(p,\mathcal M)\leq16\epsilon$. For
$\epsilon\geq1/17$ the trivial distance bound sixteen gives
$272\epsilon$, proving (79). In particular, the only zero-deficit
profiles are the members of $\mathcal M$.

For a measurable row $f:[0,1]\to[0,1]$ set
$p_i(f)=a^{-1}\int_{A_i}f$. Interpret each element $r\in\mathcal M$
as its associated step-function row. Since $r$ is Boolean,
\begin{equation}\tag{80}\label{eq:row-profile-distance}
 \lVert f-r\rVert_1=a\lVert p(f)-r\rVert_1.
\end{equation}
The pentagon link of a new vertex with row $f$ in the perfect
background $W_0$ is exactly $a^4c(p(f))$, and its degree is
$a\sum_i p_i(f)$. Indeed, deleting that vertex leaves an induced
$P_4$. An induced $P_4$ in an independent-class blow-up uses four
distinct classes, because $P_4$ has no false twins. Its four row
factors can therefore be integrated separately. The two orientations
of each such path cancel the factor $1/2$ in the rooted pentagon link.

\subsection{Continuity and variations}

The following estimates hold also on a measurable subspace of total
mass at most one. The reference values are
\[
 E_0=40a^2=\frac5{32},\qquad
 P_0=192a^5=\frac3{16384},\qquad
 \kappa_0=\frac{5P_0}{2E_0}=12a^3.
\]
For graphons $U,V$, telescoping the ten factors in the induced
pentagon integrand gives
\begin{equation}\tag{81}\label{eq:l1-continuity}
 |P(U)-P(V)|\leq\lVert U-V\rVert_1,\qquad
 |E(U)-E(V)|\leq\tfrac12\lVert U-V\rVert_1.
\end{equation}
Thus $\kappa(W)=5P(W)/(2E(W))$ is Lipschitz in a fixed neighborhood
of $W_0$.

We also allow a measurable density $z\in[0,1]$ on the vertex measure.
Write $P(W,z)$, $E(W,z)$ for the resulting integrals,
$d_z(x)=\int W(x,y)z(y)\,dy$, and $L_z(x)$ for the first variation
of $P$ at a vertex. Explicitly $L_z$ is one half the integral of the
rooted induced pentagon product over the other four vertices, with
their densities included. For a differentiable density path,
\begin{equation}\tag{82}\label{eq:vertex-density-variation}
 \frac{d}{dt}F(W,z_t)
 =E(W,z_t)^{-5/2}\int z_t'(x)
       [L_{z_t}(x)-\kappa(W,z_t)d_{z_t}(x)]\,dx.
\end{equation}
If $z$ differs from one only on a set of measure $\mu$, uniformly in
$x$ one has
\begin{equation}\tag{83}\label{eq:density-deletion-errors}
 |L_z(x)-L_1(x)|\leq2\mu,\qquad
 |d_z(x)-d_1(x)|\leq\mu.
\end{equation}
The changes in $P$ and $E$ are at most $\mu/2$ and $\mu$,
respectively. These estimates follow by telescoping the four or five
vertex density factors. Keeping the row $f=W(x,\cdot)$ fixed, but
replacing the background among the other four vertices by $W_0$, gives
\begin{equation}\tag{84}\label{eq:background-replacement-error}
 |L_1(x)-a^4c(p(f))|\leq3\lVert W-W_0\rVert_1.
\end{equation}
There are six background factors, each with the rooted prefactor $1/2$.

We need a small-degree bound that does not assume triangle-freeness:
\begin{equation}\tag{85}\label{eq:local-small-degree-link}
 L_z(x)\leq\frac14d_z(x)^2E(W,z).
\end{equation}
To see this, fix the middle edge $bc$ of the rooted pentagon
$x-a-b-c-d-x$. Drop the chord condition between $a,d$. The two
remaining neighbor choices have masses
\begin{align*}
 r&=\int W(x,u)W(u,b)(1-W(u,c))z(u)\,du,\\
 s&=\int W(x,u)W(u,c)(1-W(u,b))z(u)\,du.
\end{align*}
The inequality $u(1-v)+v(1-u)\leq1$ for $u,v\in[0,1]$ gives
$r+s\leq d_z(x)$ and $rs\leq d_z(x)^2/4$. Drop the two factors
involving the chords $xb,xc$ and integrate the middle edge, retaining
its factor $1/2$. This proves (85).

For a symmetric edge variation $H$,
\begin{equation}\tag{86}\label{eq:edge-variation}
 \begin{aligned}
 DP(W)[H]&=\frac12\int q_W(x,y)H(x,y)\,dx\,dy,\\
 DF(W)[H]&=\frac1{2E(W)^{5/2}}\int(q_W-\kappa(W))H.
 \end{aligned}
\end{equation}
Here $q_W$ is the difference of two three-variable integrals: remove
one cycle-edge factor from the induced pentagon product for the
positive term, and remove one chord factor for the negative term.
The five cycle edges and five chords, together with the factor $1/10$,
give the prefactor $1/2$.

For a Boolean step graphon $T$ put
$e_x=\int|W(x,z)-T(x,z)|\,dz$ and $e=\lVert W-T\rVert_1$. Each
of the two integrands defining $q$ has three factors incident to $x$,
three incident to $y$, and three among the other vertices.
Telescoping yields
\begin{equation}\tag{87}\label{eq:edge-variation-error}
 |q_W(x,y)-q_T(x,y)|\leq6(e_x+e_y+e).
\end{equation}
By Lemma~A.3,
\begin{equation}\tag{88}\label{eq:edge-sign-margin}
 q_{W_0}(x,y)-\kappa_0
 =\begin{cases}
   12a^3,&\text{adjacent types},\\
   -18a^3,&\text{distinct nonadjacent types},\\
   -12a^3,&\text{equal types}.
  \end{cases}
\end{equation}
The pair coefficients are multiplied by $a^3$, the mass of each old
triple. Set $g=12a^3>0$. By continuity, a Clebsch step graphon with
class masses sufficiently close to $a$ retains these signs with margin
at least $3g/4$.

\subsection{Proof of the local theorem}

\begin{proof}[Proof of Theorem A.1]
Choose a fixed $\eta>0$ small enough that the $\eta$-balls around
the seventeen reference rows are disjoint and
\[
 \frac{\eta E_0}{2}<\frac{\kappa_0}{4},\qquad
 12\eta<\frac g8.
\]
We subsequently choose $\delta_0$ sufficiently small, with $\eta$
held fixed. Write $\delta=\lVert W-W_0\rVert_1$.

\emph{Remove rows with a definite link deficit.} Let $B$ consist of
points whose rows have distance at least $\eta$ from every reference
row. Since the original template row is one of these rows,
\[
 \mu_B:=|B|\leq\delta/\eta.
\]
By (79) and (80), the perfect-background deficit on $B$ is at least
$a^4\eta/17$. Along the deletion path $z_t=1-t\mathbf1_B$,
equations (81)--(84) therefore give, uniformly for almost every
$x\in B$,
\[
 L_{z_t}(x)-\kappa(W,z_t)d_{z_t}(x)
 \leq-\frac{a^4\eta}{17}+O(\delta+\mu_B)
 \leq-\frac{a^4\eta}{34}.
\]
The implicit constant depends only on the fixed template and the
fixed $\eta$. The last inequality follows by reducing $\delta_0$.
The edge mass stays bounded away from zero. By (82), deleting $B$
increases $F$ by at least $c_B\mu_B$ for some $c_B>0$.

\emph{Remove near-zero rows.} Among the remaining points let $Z$
consist of rows whose original distance from zero is less than $\eta$.
Their original degrees are at most $\eta$, and
\[
 \mu_Z:=|Z|\leq\frac\delta{5a-\eta},
\]
because each original template row has degree $5a$. Delete $Z$ after
$B$. During this deletion its degrees are still at most $\eta$, while
$E\leq2E_0$ and $\kappa\geq\kappa_0/2$ for sufficiently small
$\delta_0$. By (85),
\[
 L_{z_t}(x)-\kappa(W,z_t)d_{z_t}(x)
 \leq-\frac{\kappa_0}{4}d_{z_t}(x)\qquad(x\in Z).
\]
If $D_Z$ denotes the edge mass removed in this step, then
\[
 D_Z=\int_0^1\int_Z d_{z_t}(x)\,dx\,dt.
\]
Thus deleting $Z$ increases $F$ by at least $c_ZD_Z$ for a constant
$c_Z>0$. The deletion of a measurable set is equivalently its
replacement by isolates.

\emph{Assign the remaining rows to clone types.} Let
$M=[0,1]\setminus(B\cup Z)$. Every point of $M$ is within $\eta$ of
a unique nonzero reference row. Call its type $\sigma(x)$; let
$\tau(x)$ be its original part label, and put
$M_i=\{x\in M:\sigma(x)=i\}$. Distinct Clebsch rows have distance
at least $6a$, so the measure $\nu$ of the misassigned set
$\{\sigma\ne\tau\}$ satisfies
\[
 \nu\leq\frac\delta{6a-\eta}.
\]
In particular, for $s_i=|M_i|$,
\begin{equation}\tag{89}\label{eq:reclassified-masses}
 |s_i-a|\leq\mu_B+\mu_Z+\nu=O(\delta).
\end{equation}
On $M$ let $T(x,y)=A_{\sigma(x),\sigma(y)}$ be the ideal Clebsch
step graphon. For almost every $x\in M$,
\begin{equation}\tag{90}\label{eq:reclassified-errors}
 \int_M|W(x,y)-T(x,y)|\,dy\leq\eta+\nu,\qquad
 e_M:=\lVert W|_{M^2}-T\rVert_1\leq\delta+2\nu.
\end{equation}
For the first inequality, the assigned reference row differs from its
reclassified version only on the misassigned set. For the second, $T$
differs from the original template on pairs having at least one
misassigned endpoint.

\emph{Repair all remaining edge errors.} Consider
$W_t=T+t(W|_{M^2}-T)$ for $0\leq t\leq1$, on the measure space $M$.
The row errors and global error remain bounded by the right-hand
sides of (90) throughout this path. By (89), the margin for $T$ is at
least $3g/4$. Equations (81), (87) and (90) imply
\[
 |(q_{W_t}-\kappa(W_t))-(q_T-\kappa(T))|
 \leq12\eta+O(\delta)<g/4.
\]
Hence $q_{W_t}-\kappa(W_t)\geq g/2$ where $T=1$, and it is at most
$-g/2$ where $T=0$. The variation $H=W|_{M^2}-T$ has the opposite
sign on each region. From (86),
\[
 \frac{d}{dt}F(W_t)\leq-c_Ee_M
\]
for a fixed $c_E>0$. Integrating gives
$F(T)-F(W|_{M^2})\geq c_Ee_M$.

\emph{Balance the classes and estimate the distance.} Put $q=|M|$
and $V=(\sum_i(s_i-q/16)^2)^{1/2}$. The Hessian estimate (78) applies
by (89). Combining the four comparisons gives
\begin{equation}\tag{91}\label{eq:four-defects}
 F(W_0)-F(W)\geq c_B\mu_B+c_ZD_Z+c_Ee_M+c_HV^2.
\end{equation}
Isolating $B$ costs at most $2\mu_B$ in $L^1$, and then isolating $Z$
costs exactly $2D_Z$. Replacing the remaining graphon by $T$ costs
$e_M$. Within $M$, balance the class masses by moving a set of measure
$\tfrac12\sum_i|s_i-q/16|$ between classes; atomlessness permits this
operation. The resulting graphon belongs to $\mathcal E_{\mathrm{Cl}}$,
and this last change costs at most $\sum_i|s_i-q/16|\leq4V$. Thus
\[
 d_1(W,\mathcal E_{\mathrm{Cl}})\leq2\mu_B+2D_Z+e_M+4V.
\]
Since $\mu_B,D_Z,e_M\leq1$, its square is at most
$16\mu_B+16D_Z+4e_M+64V^2$. Together with (91) this proves (76),
after multiplying by $5^{5/2}$ and adjusting the positive constant.

If $R_{\mathrm{ind}}(W)=R_\ast$, all four nonnegative terms in (91)
vanish. Thus $B$ is null, $Z$ consists of actual isolates, the
remaining graphon is exactly $T$, and its sixteen class masses are
equal. This proves the equality assertion without assuming that the
distance infimum is attained. The converse follows by homogeneity.
\end{proof}
